\begin{center}
{\Large \bfseries Almacenamiento y Análisis Eficiente de Resultados de Simulaciones Hidráulicas mediante TileDB}

\vspace{1cm}
{\Large \bfseries RESUMEN}

\vspace{1.5cm}
\end{center}


Este trabajo diseña e implementa un sistema de almacenamiento y análisis eficiente para simulaciones hidráulicas de alta resolución. Las simulaciones del simulador Triton generan hasta 235 GB de datos GeoTIFF por escenario; el \textit{pipeline} ETL desarrollado los transforma en un array TileDB sparse con filtro ByteShuffle+ZSTD-9, reduciéndolos a 17,5 GB (reducción del 92,7\,\%, ratio 13,7$\times$) al almacenar exclusivamente las celdas húmedas, que representan solo el $\sim$8,3\,\% del dominio.

El sistema incluye un motor de 24 consultas analíticas en ocho bloques, que cubren desde consultas básicas de calado y velocidad hasta indicadores de inestabilidad hidrodinámica según los criterios de Russo et al.\ y Xia et al., zona de graves daños normativa (Real Decreto 9/2008) y tiempo de evacuación disponible. Un \textit{benchmark} frente a GeoTIFF muestra ventajas de hasta 12,4$\times$ en cinco de seis patrones de acceso, con resultados consistentes sobre cuatro \textit{datasets} independientes. La validación funcional confirma la exactitud de las consultas mediante comparación sistemática con los datos originales.

El sistema se ha aplicado a 56 simulaciones sobre una cuenca de $\sim$78\,000 km$^2$ en el centro de Estados Unidos, y es generalizable a las salidas de cualquier simulador hidráulico 2D con estructura equivalente. La interfaz de usuario es una aplicación web Streamlit con mapas interactivos, comparativa multi-escenario y exportación a GeoTIFF, desplegada en producción sobre un servidor remoto.


